% foundations/pose3_on_ellipsoid.tex

\chapter{Spatial Pose in the Earth Realization}

The preceding chapter produced the realized surface trajectory,
\(\mathbf t\), \(\mathbf n_{\mathcal E}\), \(\mathbf b\), and
\(\kappa_g\). These are not a second alignment definition. They are the
metric outputs required to ask a spatial engineering question: how do
cant, gravity, and speed act at each position of the same alignment?

\section{From surface trajectory to pose3}

Applying the frame operator gives the uncanted surface frame
\[
F_{\mathcal E}(s)=
\bigl(\mathbf t(s),\mathbf b(s),\mathbf n_{\mathcal E}(s)\bigr).
\]
The surface-based pose is consequently
\[
\mathrm{pose3}_{\mathcal E}(s)=
\bigl(
\gamma_{\mathcal E}(s),
\mathbf t(s),
\mathbf n_{\mathcal E}(s),
\kappa_g(s),
\phi(s)
\bigr).
\]
Planar heading \(\theta\) has not been discarded arbitrarily: its role
is now carried by a tangent vector field whose comparison from point to
point depends on the surface connection.

The curvature edit changes \(\kappa_g\) directly in this realization and
therefore changes the propagated tangent, surface trajectory, and local
frame. The reference ellipsoid and gravity model remain fixed context;
cant remains fixed unless edited or coupled by an applicable constraint.

\section{Cant rotates the realized cross-section}

Cant rotation \(\phi(s)\) acts about the tangent:
\[
\mathbf b_{\phi}(s)
=R_{\mathbf t}(\phi(s))\,\mathbf b(s),
\qquad
\mathbf n_{\phi}(s)
=R_{\mathbf t}(\phi(s))\,\mathbf n_{\mathcal E}(s).
\]
This operation receives the realized surface frame and cant state and
produces the canted cross-section. It preserves alignment identity and
retains both state and realization provenance. It does not infer cant
from a rendered rail model or decide that the resulting balance is
acceptable.

\section{Gravity makes the pose physically interpretable}

Let \(\mathbf g(s)\) be the applicable local gravity vector at
\(\gamma_{\mathcal E}(s)\). A simplified realization may take
\(\mathbf g\parallel-\mathbf n_{\mathcal E}\); a higher-fidelity
realization may distinguish ellipsoid normal, geoid normal, and local
gravity. The choice belongs to the metric and physical evaluation
context, not to alignment identity.

Gravity is introduced here because cant has an engineering consequence
only relative to a force direction. The surface normal constructs the
frame; the gravity vector supplies the physical direction used by the
evaluation.

\section{Curvature, speed, cant, and effective acceleration}

For speed \(v\), the curvature-induced lateral term is
\[
a_{\kappa}(s)=v^2\kappa_g(s).
\]
Projecting the curvature and gravity contributions into the canted
lateral direction gives, under the adopted sign convention,
\[
a_{\mathrm{eff}}(s)
=v^2\kappa_g(s)
-\mathbf g(s)\cdot\mathbf b_{\phi}(s).
\]
Under the local planar approximation this reduces to
\[
a_{\mathrm{eff}}(s)
=v^2\kappa(s)-g\sin\phi(s).
\]

The dependency is now explicit: the curvature edit changes
\(\kappa_g\) and the realized frame; speed supplies the operational
condition; cant rotates the cross-section; gravity supplies the physical
direction; evaluation produces effective acceleration. The result
depends on all five inputs. It cannot decide an admissible speed, approve
the edit, or replace the engineering rule under which it is assessed.

\section{What this specialization establishes}

\begin{description}
\item[Usable now:] AIM can preserve identity and provenance while
  reconstructing an intrinsic state, applying an explicit metric
  context, producing planar representations, and evaluating qualified
  consequences. A local planar realization remains sufficient whenever
  its declared accuracy is sufficient for the task.
\item[Formal specialization:] These chapters state how the same operator
  contracts specialize to a reference ellipsoid, including the
  geodesic/normal decomposition, surface pose, canted frame, and
  gravity-aware dependency. This is a mathematical specialization, not
  a claim of completed production software.
\item[Research:] Validated applicability criteria, error bounds
  between planar/projected and ellipsoidal realizations, robust numerical
  realization over project-scale data, and the required fidelity of
  gravity models remain research and validation work.
\item[Practical scope:] AIM does not require every edit to
  run on an ellipsoid, replace established project CRS workflows, or use
  a geodesic solver when a qualified planar realization answers the
  engineering question.
\end{description}

Earth curvature has therefore earned a limited and exact place in the
formal machinery. It becomes necessary when the realization context
makes surface effects material; otherwise it remains an available
specialization. In both cases the alignment identity, constructive
curvature history, and responsibility boundaries remain intact.
